% Integrated after independent coordinating-agent audit.
\begin{claim}{C.24}{Refuted}{Normalization propagation from scalar adjoints alone}\label{C24}
Completeness, cocompleteness, enough projectives, an exact scalar-copower embedding with both adjoints, and the canonical normalization $L\varphi(\mathbf1)\simeq k$ do not by themselves imply
\[
 L\varphi\kappa(V)\longrightarrow V
 \quad\text{is invertible for every free scalar module }V.
\]
We give a counterexample over a field, with a specified object $\mathbf1$ and its actual scalar copower. No symmetric monoidal tensor with that object as unit has been constructed. Accordingly this refutes a proposed argument omitting the monoidal data, not the manuscript's theorem or its full hypotheses.
\end{claim}
\begin{proof}
Fix a field $k$ and define the additive $k$-linear functor
\[
 E(V)=\operatorname{coker}(V\longrightarrow V^{**}),
 \qquad V^*=\operatorname{Hom}_k(V,k),
\]
using the evaluation map. Vector-space duality is exact, and evaluation is injective. Applying the kernel--cokernel sequence to the natural map from a short exact sequence to its double dual proves that $E$ is exact. It vanishes on finite-dimensional vector spaces.

It does not vanish on an infinite free vector space $V=k^{(I)}$. Indeed $V^*=k^I$. We choose a nonzero linear functional on the nonzero quotient $k^I/k^{(I)}$ and compose with the quotient map. This gives a nonzero element of $V^{**}$ vanishing on every finitely supported vector in $k^I$. An evaluation functional coming from $V$ is a finite linear combination of coordinate functionals; if it vanishes on all finitely supported vectors, every coefficient is zero. Thus our chosen functional has a nonzero class in $E(V)$. No cardinality bound or large-cardinal assumption enters this argument.

On $\mathrm{Vect}_k\times\mathrm{Vect}_k$ put $H(U,V)=(E(V),U)$ and use the square-zero monad $T=\mathrm{Id}\oplus H$. Its unit includes the identity summand, and its multiplication, under
\[
 T^2=\mathrm{Id}\oplus H\oplus H\oplus H^2,
\]
sends $(x,h_1,h_2,h_{12})$ to $(x,h_1+h_2)$. Additivity verifies the monad identities. Its algebra category $\mathcal A$ has objects
\[
 (U,V,a:U\to V,b:E(V)\to U),
 \qquad ab=0,\quad bE(a)=0.
\]
A morphism $(f,g)$ satisfies $ga=a'f$ and $fb=b'E(g)$.

We verify the categorical hypotheses explicitly. Kernels and cokernels are computed on the two underlying vector spaces, since exactness of $E$ supplies the induced maps and relations. Finite sums are likewise componentwise. The coimage-to-image map is an isomorphism because it is so on the underlying vector spaces and the forgetful functor reflects isomorphisms. Hence $\mathcal A$ is abelian, with an exact faithful forgetful functor $U_{\mathcal A}$ to $\mathrm{Vect}_k^2$.

All small limits exist and are created by $U_{\mathcal A}$. For example a product has underlying pair $(\prod_iU_i,\prod_iV_i)$; its $b$ is induced by
\[
 E(\prod_iV_i)\longrightarrow E(V_i)\xrightarrow{b_i}U_i
\]
for all $i$. This does not assume that $E$ preserves products. The free algebra functor $F$ is left adjoint to $U_{\mathcal A}$, with underlying functor $T$. Every free algebra is projective, because
\[
 \operatorname{Hom}_{\mathcal A}(F(W),-)
 =\operatorname{Hom}_{\mathrm{Vect}_k^2}(W,U_{\mathcal A}(-))
\]
is exact. The action map $F(U_{\mathcal A}X)\to X$ is epic, since it has an underlying vector-space section given by the monad unit. Thus $\mathcal A$ has enough projectives.

We also check all small coproducts, without assuming accessibility of $T$. Each algebra $X$ is the coequalizer of its standard free presentation
\[
 F(TU_{\mathcal A}X)\rightrightarrows F(U_{\mathcal A}X)\longrightarrow X.
\]
For a set of algebras $X_i$, take free algebras on the corresponding direct sums of underlying vector spaces and take the cokernel of the difference of the two induced maps. The free adjunction and the displayed coequalizer property identify maps from this cokernel with arbitrary families of maps from the $X_i$. It is therefore their coproduct. Abelian cokernels then give all small colimits. In particular we are not claiming that these coproducts are componentwise sums.

Now define
\[
 \kappa(W)=(W,0,0,0),\qquad
 \varphi(U,V,a,b)=\operatorname{coker}b,\qquad
 \psi(U,V,a,b)=\ker a.
\]
The morphism equations give actual adjunctions
\[
 \varphi\dashv\kappa\dashv\psi.
\]
Indeed, a map from $(U,V,a,b)$ to $\kappa(W)$ is a map $U\to W$ annihilating $b$, whereas a map from $\kappa(W)$ to $(U,V,a,b)$ is a map $W\to U$ annihilated by $a$. The embedding $\kappa$ is fully faithful and exact. It preserves all limits, because those limits on its image are componentwise, and all colimits, because it is left adjoint to $\psi$. For the specified object $\mathbf1=\kappa(k)$, the latter adjunction is precisely the scalar-copower formula
\[
 \operatorname{Hom}_{\mathcal A}(\kappa(W),X)
 \cong\operatorname{Hom}_k\bigl(W,\operatorname{Hom}_{\mathcal A}(\mathbf1,X)\bigr).
\]
Thus $\kappa$ is the actual copower of $\mathbf1$, not a replacement functor bearing the same name.

The free algebras on the two sorts are
\[
 P_0(W)=(W,W,1,0),\qquad
 P_1(W)=(E(W),W,0,1).
\]
Both are projective for every $W$. Writing $S_1(W)=(0,W,0,0)$, we have short exact sequences
\[
 \begin{aligned}
 0&\longrightarrow S_1(W)\longrightarrow P_0(W)
      \longrightarrow\kappa(W)\longrightarrow0,\\
 0&\longrightarrow\kappa(E(W))\longrightarrow P_1(W)
      \longrightarrow S_1(W)\longrightarrow0.
 \end{aligned}
\]
The first epimorphism has component maps $(1,0)$ and the second $(0,1)$. Iterating gives a projective resolution of $\kappa(W)$ whose terms in homological degrees $2n$ and $2n+1$ are, respectively, $P_0(E^nW)$ and $P_1(E^nW)$. The odd differential has components $(0,1)$ and each positive even differential has components $(1,0)$. The two displayed short exact sequences verify exactness throughout.

Applying $\varphi$ gives $E^nW$ in degree $2n$ and zero in every odd degree, with zero differentials. Therefore
\[
 L_{2n}\varphi(\kappa W)=E^nW,
 \qquad L_{2n+1}\varphi(\kappa W)=0.
\]
The degree-zero counit is canonically $1_W$. For $W=k$, all positive groups vanish because $E(k)=0$; hence the canonical normalization $L\varphi(\mathbf1)\to k$ is a quasi-isomorphism. For an infinite free $W$, however, $L_2\varphi(\kappa W)=E(W)\ne0$, and its derived counit is not a quasi-isomorphism.

Finally we can see the failure of AB4 directly. Since $E(k)=0$, $S_1(k)=P_1(k)$, so $P_1(k)\to P_0(k)$ is monic. For $W=k^{(I)}$, the coproduct of these maps is $P_1(W)\to P_0(W)$, because both families are free algebras. Its kernel is $\kappa(E(W))\ne0$. Thus even an infinite coproduct of copies of this one monomorphism need not be monic. The projective generator $P_0(k)\oplus P_1(k)$ is not small: its Hom functor fails to preserve these coproducts. This does not satisfy the corollary's finitely presented projective-generator hypothesis. Most importantly, we have supplied no tensor or interval satisfying the remaining original assumptions, so the example is not being used to refute those assumptions' full conclusion.
\end{proof}
